\chapter{实验设计}

\section{测评实验设置}

本章实验围绕“同一批文档问题在纯文本 RAG 与多模态 RAG 下的效果差异”展开。评测统一采用 DocVQA 和 ChartQA 两类公开任务：DocVQA 侧重文档页面中的文字、表格、票据、广告和表单信息抽取；ChartQA 侧重图表数值读取、趋势判断和视觉区域理解。两类任务分别覆盖文本密集型文档问答与视觉密集型图表问答，能够同时检验系统的文档解析、检索重排、模型选择、局部视觉证据组织和多模态生成能力。

实验使用同一组问题和文档分别运行 Text-RAG 与 Our-RAG，并使用数据集自带的参考答案字段进行统一评分。Text-RAG 只基于解析得到的文本 chunk 完成检索和生成，Our-RAG 在此基础上进一步引入页面图像、局部 region crop、图表提示和 VLM 视觉理解。所有样本均记录预测答案、命中证据、引用页码、视觉区域、模型选择结果、延迟和错误标签，并通过同一套脚本完成短答案抽取与指标计算。这样可以保证不同方法在输入样本、评分逻辑和运行记录上保持一致，后续结果差异主要反映多模态证据、检索重排与 GRPO 决策机制带来的影响。

\section{数据集}

本文选择 DocVQA 与 ChartQA 两类数据作为评测来源：

\begin{itemize}
  \item \textbf{DocVQA}：考察文档页面、表格、广告、表单和票据等场景中的文字定位、版面理解和答案抽取能力。
  \item \textbf{ChartQA}：考察柱状图、折线图、数值比较和视觉趋势理解能力。
  \item \textbf{覆盖问题类型}：包括文档文字抽取、表格行列定位、广告与商品标签识别、票据版面理解、图表数值读取和趋势判断。
  \item \textbf{数据清洗}：评测前剔除标注歧义或无法稳定复现的样本，并保留数据集原始答案字段、页面来源和问题类型信息。
\end{itemize}

\section{对比方法}

本文比较两类方法：一类是主结果中的完整模型对比，另一类是针对系统关键组件的消融实验。

\begin{itemize}
  \item \textbf{Text-RAG}：只使用文档文本抽取结果构建 chunk 和向量索引。该方法适合章节总结、定义、普通文本问答等场景，是多模态系统的基础 baseline。
  \item \textbf{Our-RAG}：在 Text-RAG 基础上加入页面视觉 evidence、局部高清 crop、图表提示、Reranker、统一多模态向量化、GRPO 决策模型和 VLM 推理。该方法能处理图表、视觉区域、广告品牌、商品标签、表格和图像型问题。
  \item \textbf{Our-RAG with Rule Decision}：将 GRPO 决策模型替换为硬编码规则，用关键词和问题类型规则选择文本 LLM 或多模态 VLM，用于观察决策模型的贡献。
  \item \textbf{Our-RAG without Reranker}：去掉 BGE-Reranker-v2-m3，只使用向量召回结果进入生成阶段，用于观察二次重排对证据质量的影响。
  \item \textbf{Our-RAG without Unified MM Embedding}：去掉统一文本-图像向量化接口，仅使用文本向量召回视觉区域关联描述，用于观察多模态向量融合对召回和生成的影响。
\end{itemize}

所有方法共享同一套短答案抽取和评分脚本，以保证比较公平。对于生成式回答，系统优先抽取 Final answer: 后的短答案，并进行大小写、标点和空格等基础格式处理，避免解释性长句直接干扰短答案指标。

\section{评测指标}

本文使用以下指标：

\begin{itemize}
  \item \textbf{Exact Match (EM)}：预测短答案与参考答案严格匹配的比例。该指标适合衡量抽取式或短答案问答。
  \item \textbf{ANLS}：平均归一化 Levenshtein 相似度，是 DocVQA 常用指标，能够容忍轻微拼写或格式差异。
  \item \textbf{Answer Match}：宽松答案命中率，用于衡量模型输出是否与参考答案形成有效匹配，可辅助观察生成式回答在格式差异下的实际可用性。
  \item \textbf{Latency}：单个问题平均推理耗时，用于评估系统可用性。
\end{itemize}

此外，实验记录中也保留 Recall@K、MRR 和 Citation Accuracy，用于检查检索召回、证据排序和 citation 对齐是否稳定。最终结果表重点展示 EM、ANLS、Answer Match 和平均延迟；检索相关指标作为辅助分析记录，用于定位证据召回和溯源对齐问题。

\section{运行环境}

本项目实验在本地工作站上完成，使用的 GPU 配置如下：

\begin{itemize}
  \item NVIDIA GeForce RTX 5090 32GB $\times$ 1
  \item NVIDIA GeForce RTX 3070 8GB $\times$ 1
\end{itemize}

\section{实验流程}

实验流程分为四步。第一步，使用统一脚本读取 DocVQA 和 ChartQA 样本，并构建文档对象。第二步，分别运行 Text-RAG、Our-RAG 及消融配置，保存每个样本的预测答案、top evidence、引用页、模型选择结果、延迟和失败标签。第三步，对模型输出进行短答案抽取，计算 EM、ANLS 和 Answer Match。第四步，对错误样本进行人工检查，分析错误来源是否来自检索、视觉区域选择、OCR、VLM 推理或答案格式波动。

针对不同类型问题，系统在运行阶段采用相同的检索、重排和生成流程。对于 DocVQA，系统重点利用页面 OCR、局部视觉区域和版面信息；对于 ChartQA，系统优先生成 chart region 并增强图表任务提示，使 VLM 能够围绕坐标轴、图例、柱高、折线变化和数值比较关系进行回答。
